% ----------------------------------------------------------
% Fundamentação Teórica
% ----------------------------------------------------------
\chapter{Fundamentação Teórica} \label{cha:fundamentacao}

Este capítulo explica as ideias que servem de base para o projeto: como a gestão de conteúdo digital surgiu e mudou ao longo do tempo, como ela funciona hoje e quais teorias ajudam a entender o problema.

\section{Histórico do Assunto}

A gestão de conteúdo digital vem dos arquivos de papel e mudou bastante com o surgimento dos \textit{Content Management Systems} (CMS), na década de 1990. Segundo Boiko \cite{boiko2004}, um CMS é uma aplicação \textit{web} que permite que pessoas sem formação técnica criem, atualizem e publiquem conteúdo digital. Pela primeira vez, escrever o conteúdo ficou separado de programar como ele aparece na tela. Essa separação facilitou a criação de conteúdo na \textit{web} e deu origem aos sistemas de gestão que existem hoje.

Essa mudança cresceu com a popularidade das plataformas de vídeo e dos ambientes de educação online, como Coursera e Udemy. Para Rosenfeld, Morville e Arango \cite{rosenfeld2015}, a Arquitetura da Informação serve para que os ambientes digitais sejam fáceis de encontrar, de navegar e de entender. Isso fica mais importante quando o conteúdo cresce e fica mais complexo. Por isso, ferramentas de organização feitas para a área da educação deixam de ser um diferencial e passam a ser uma necessidade.

Essa complexidade afeta principalmente quem cria conteúdo sozinho. Diferente de uma empresa, onde cada pessoa tem uma função, o criador de cursos digitais faz o papel de roteirista, produtor, editor e gestor do projeto ao mesmo tempo. Como diz Drucker \cite{drucker1999}, o trabalhador do conhecimento decide sozinho qual é a sua tarefa. Por isso, ele precisa de liberdade, mas também de organização, para não perder o foco. Sem um sistema que siga a estrutura natural de um curso (em que aulas dependem de roteiros, roteiros dependem de pesquisas e grupos de aulas formam módulos), o trabalho se espalha entre aplicativos de notas, planilhas e pastas de arquivos sem padrão.

\section{Atualidade do Assunto}

Hoje, muitas empresas usam o que se chama de \textit{Content Operations} (ContentOps). Segundo o Content Marketing Institute \cite{mcknight2023}, o termo se refere a tudo o que envolve conteúdo dentro de uma organização: o planejamento, a criação, o controle de qualidade, a medição de resultados e a organização do material. A mesma fonte explica que a ideia é que todos produzam conteúdo seguindo um mesmo jeito de trabalhar, em vez de cada pessoa ter o seu próprio processo.

Na parte técnica, esse movimento usa sistemas \textit{headless}, que separam a criação do conteúdo da forma como ele aparece na tela. Assim, o conteúdo é feito uma única vez e pode ser publicado em vários canais, como site, aplicativo e redes sociais. Esse modelo é chamado de \textit{omnichannel}. O Gartner previu que, até 2026, pelo menos 70\% das organizações precisariam adotar plataformas de experiência digital montadas em partes \cite{demeny2025}.

Para quem cria conteúdo sozinho, porém, essas práticas ainda são pouco comuns. As soluções existentes foram feitas para empresas grandes, e faltam ferramentas pensadas para quem trabalha sozinho ou em equipes pequenas. É essa lacuna que o Creator Workbench quer preencher.

\section{Outros Contextos do Assunto}

Dois campos de estudo ajudam a entender o problema deste trabalho: a Andragogia e a Arquitetura da Informação.

\subsection{Andragogia}

A Andragogia é definida por Malcolm Knowles, seu principal autor, como a arte e a ciência de ensinar adultos. Na obra \textit{The Modern Practice of Adult Education: Pedagogy vs. Andragogy} \cite{knowles1970}, Knowles defende que a educação de adultos precisa ser diferente da educação de crianças. Entre os princípios que ele aponta estão a necessidade de saber por que está aprendendo, a autonomia de quem aprende, o uso da experiência anterior e o foco em resolver problemas.

Nos cursos digitais, essa teoria ajuda a pensar no conteúdo. Um curso organizado de forma clara, em ordem e com ligações visíveis entre suas partes ajuda o aluno a aprender com mais autonomia e a guardar melhor o que aprendeu. Ao organizar módulos e aulas de forma nativa, o Creator Workbench dá ao criador as condições técnicas para aplicar esses princípios no seu curso.

Em \textit{The Adult Learner} \cite{knowlesadult}, publicada no Brasil como \textit{Aprendizagem de Resultados}, Knowles apresenta os seis princípios da aprendizagem de adultos. A obra é um clássico da educação em empresas e traz pontos que devem ser considerados em qualquer ação educacional.

Entre esses princípios, dois combinam bem com este projeto: aprender para resolver problemas e ver utilidade imediata no conteúdo. Um criador que não consegue organizar o próprio trabalho dificilmente vai montar um curso que respeite esses princípios para os seus alunos.

\subsection{Arquitetura da Informação}

A Arquitetura da Informação ajuda a pensar como organizar os sistemas de criação. Em \textit{Information Architecture: For the Web and Beyond} \cite{rosenfeld2015}, Rosenfeld, Morville e Arango explicam que a área busca tornar os ambientes digitais fáceis de encontrar, de navegar e de entender. Seus conceitos principais são organização, rotulação, navegação e busca.

No trabalho com conteúdo, essa ideia mostra por que ferramentas genéricas, como Trello e Notion, muitas vezes não funcionam bem para cursos: elas não têm uma estrutura própria para a educação. Um \textit{board} de \textit{Kanban} não diferencia uma ``aula'' de uma ``tarefa administrativa''. Um banco de dados do Notion, por padrão, não tem a ideia de ``dependência entre roteiros''. O Creator Workbench coloca essa estrutura dentro do próprio sistema e aplica as ideias de Rosenfeld et al. à criação de cursos digitais. Assim, o sistema fica organizado e também fácil de entender para quem o usa.